\section{Introduction}
\label{sec:introduction}

Fraud in online platforms---ranging from fake product reviews on Yelp to money laundering on the Bitcoin blockchain---causes billions of dollars in annual losses and erodes user trust in digital ecosystems~\cite{wang2019fdgars}. Detecting fraudulent behavior is fundamentally challenging because fraudsters continuously adapt their strategies to evade detection, and fraudulent instances typically constitute a tiny minority of all transactions, leading to severe class imbalance.

Graph-based approaches have become the dominant paradigm for fraud detection because they naturally model the relational structure of fraudulent activity: fraudsters interact with victims, form collusive rings, and leave topological traces that are invisible to non-relational models~\cite{liu2020alleviating,dou2020enhancing}. Graph Neural Networks (GNNs)~\cite{kipf2016semi,velickovic2017graph,hamilton2017inductive} learn node representations by iteratively aggregating features from local neighborhoods, enabling the model to capture both attribute-based signals and structural context. However, standard GNNs face two well-documented challenges in fraud settings: (i) camouflage, where fraudsters mimic benign feature distributions to blend in~\cite{dou2020enhancing}, and (ii) neighborhood inconsistency, where a fraudster's neighbors are predominantly benign nodes, causing message passing to dilute the fraudulent signal~\cite{liu2020alleviating}.

Several GNN-based fraud detectors have been proposed to address these issues. CARE-GNN~\cite{dou2020enhancing} uses a reinforcement-learning-based neighbor selector to filter out suspicious relations. PC-GNN~\cite{liu2021pick} applies a label-balanced neighborhood sampler to mitigate class imbalance during training. GraphConsis~\cite{liu2020alleviating} explicitly models context, feature, and relation inconsistencies. While effective, these methods modify the GNN architecture or the message-passing procedure, adding complexity and dataset-specific design choices.

More recently, GAAP~\cite{duan2025global} introduced the concept of attribute-association patterns: by adaptively binning continuous features and combining them with neighborhood information via message passing, GAAP learns pattern representations that distinguish fraudulent from benign behaviors. A global aggregation module then captures correlations across these patterns, enabling the model to recognize fraud at multiple granularities. GAAP achieves state-of-the-art results against 24 baselines on 7 datasets. Its design is notably elegant: rather than engineering complex neighbor selection or sampling strategies, GAAP shifts the focus to what the attribute-association patterns themselves reveal about fraud.

Despite its strong performance, GAAP uses a standard cross-entropy loss that assigns equal weight to every training node. In fraud detection, this is suboptimal because the incidence of any specific attribute-association pattern is highly skewed: common patterns (e.g., a reviewer with a moderate rating history, few friends, and no flagged content) dominate the training set, while rare patterns that signal genuine fraud (e.g., a reviewer with an anomalously rapid burst of five-star reviews followed by account dormancy) appear so infrequently that their contribution to the gradient is negligible. The model has no incentive to specialize on these rare but diagnostically valuable patterns.

This observation motivates our central question: can we improve GAAP's fraud detection capability simply by telling the optimizer which training examples encode rare patterns? Concretely, we hypothesize that if nodes with low-frequency attribute-association patterns receive higher weight during loss computation, the model will learn stronger decision boundaries around fraud-specific behaviors that are otherwise underrepresented in the training distribution.

We propose RP-GAAP (Rare-Pattern-weighted GAAP), a lightweight extension that applies sample-level weighting based on pattern rarity without modifying the GAAP architecture. Our contributions are:

\begin{itemize}
    \item We introduce a simple, non-parametric rarity scoring mechanism that discretizes the top-$k$ features of each node into quantile bins, counts pattern co-occurrence frequencies across the training set, and assigns higher loss weights to nodes whose attribute-association patterns are underrepresented.
    \item We integrate these rarity scores as per-sample weights into GAAP's cross-entropy loss, yielding the RP-GAAP training objective with only three hyperparameters (number of bins $b$, number of features $k$, and maximum weight $w_{\text{max}}$), requiring no modification to the model forward pass.
    \item We evaluate RP-GAAP against the original GAAP on four real-world fraud detection benchmarks (YelpChi, T-Finance, Elliptic, and Tolokers), demonstrating improved fraud recall and Recall@K on YelpChi and T-Finance---with a $+0.97$ percentage-point gain on T-Finance---alongside ablations comparing rare-pattern weighting against focal loss~\cite{lin2017focal} and their combination.
\end{itemize}

The remainder of this paper is organized as follows. Section~\ref{sec:related_work} reviews related work on GNN-based fraud detection and loss re-weighting. Section~\ref{sec:methodology} describes GAAP, our rarity quantification method, and the RP-GAAP training objective. Section~\ref{sec:experiment_setups} details the experimental setup. Section~\ref{sec:results} presents and discusses results. Section~\ref{sec:conclusion} concludes and outlines directions for future work.
